% table_sota.tex
% Table II — Comparison with state-of-the-art compression/quantization methods
%
% Key references (verified PPL from official repos/papers):
%   [GPTQ]       Frantar et al., ICLR 2023. OPT-125M WikiText-2 4-bit PPL = 31.12
%   [SqueezeLLM] Kim et al., ICML 2024. LLaMA-7B WikiText-2 3.45-bit PPL = 5.79
%   [LLM.int8]   Dettmers et al., NeurIPS 2022. OPT->=6.7B, near-lossless 8-bit
%
% IMPORTANT: PPL values are NOT directly comparable — different models & datasets.
% The table contrasts methodology and deployment characteristics.
%
% Usage: % table_sota.tex
% Table II — Comparison with state-of-the-art compression/quantization methods
%
% Key references (verified PPL from official repos/papers):
%   [GPTQ]       Frantar et al., ICLR 2023. OPT-125M WikiText-2 4-bit PPL = 31.12
%   [SqueezeLLM] Kim et al., ICML 2024. LLaMA-7B WikiText-2 3.45-bit PPL = 5.79
%   [LLM.int8]   Dettmers et al., NeurIPS 2022. OPT->=6.7B, near-lossless 8-bit
%
% IMPORTANT: PPL values are NOT directly comparable — different models & datasets.
% The table contrasts methodology and deployment characteristics.
%
% Usage: % table_sota.tex
% Table II — Comparison with state-of-the-art compression/quantization methods
%
% Key references (verified PPL from official repos/papers):
%   [GPTQ]       Frantar et al., ICLR 2023. OPT-125M WikiText-2 4-bit PPL = 31.12
%   [SqueezeLLM] Kim et al., ICML 2024. LLaMA-7B WikiText-2 3.45-bit PPL = 5.79
%   [LLM.int8]   Dettmers et al., NeurIPS 2022. OPT->=6.7B, near-lossless 8-bit
%
% IMPORTANT: PPL values are NOT directly comparable — different models & datasets.
% The table contrasts methodology and deployment characteristics.
%
% Usage: % table_sota.tex
% Table II — Comparison with state-of-the-art compression/quantization methods
%
% Key references (verified PPL from official repos/papers):
%   [GPTQ]       Frantar et al., ICLR 2023. OPT-125M WikiText-2 4-bit PPL = 31.12
%   [SqueezeLLM] Kim et al., ICML 2024. LLaMA-7B WikiText-2 3.45-bit PPL = 5.79
%   [LLM.int8]   Dettmers et al., NeurIPS 2022. OPT->=6.7B, near-lossless 8-bit
%
% IMPORTANT: PPL values are NOT directly comparable — different models & datasets.
% The table contrasts methodology and deployment characteristics.
%
% Usage: \input{tables/table_sota}

\begin{table*}[t]
\centering
\caption{Comparison with representative post-training compression methods.
         \textbf{Note:} PPL values are evaluated on different models and datasets
         (see Model column) and are \emph{not} directly comparable.
         The table contrasts methodological characteristics and deployment target.
         $\ddagger$ FP32 footprint of counted tensors;
         $\star$ compression vs.\ FP32;
         $\dagger$ reported on WikiText-2, not our TinyStories corpus.}
\label{tab:sota}
\setlength{\tabcolsep}{4pt}
\renewcommand{\arraystretch}{1.2}
\begin{tabular}{llccccccc}
\toprule
\textbf{Method} & \textbf{Model} & \textbf{Bits/w} & \textbf{Size} & \textbf{Ratio$^\star$} &
\textbf{PPL\,$\downarrow$} & \textbf{HW} & \textbf{Flash-} & \textbf{Meta-} \\
 & & & & & & \textbf{target} & \textbf{stream} & \textbf{heuristic} \\
\midrule
FP32 (ours)               & TinyStories-15M & 32   & 60.8\,MB$^\ddagger$ & $1\times$      & \phantom{0}6.92  & MCU     & \texttimes & \texttimes \\
\midrule
GPTQ~\cite{frantar2023gptq}       & OPT-125M        & 4    & $\sim$15\,MB    & $\sim 8\times$ & 31.12$^\dagger$  & GPU     & \texttimes & \texttimes \\
SqueezeLLM~\cite{kim2024squeezellm}& LLaMA-7B        & 3.45 & $\sim$3.5\,GB   & $\sim 4\times$ & \phantom{0}5.79$^\dagger$ & GPU & \texttimes & \texttimes \\
LLM.int8()~\cite{dettmers2022llmint8}& OPT-$\geq$6.7B& 8    & $\sim 2\times$ red. & $\sim 2\times$ & $\approx$FP16$^\dagger$ & GPU & \texttimes & \texttimes \\
\midrule
INT4 Uniform (ours)       & TinyStories-15M & 4    & 7.61\,MB        & $7.99\times$   & 11070            & \textbf{MCU} & \checkmark & \texttimes \\
INT4 Block-wise (ours)    & TinyStories-15M & $\approx$5 & 9.51\,MB  & $6.39\times$   & \phantom{0}10.71 & \textbf{MCU} & \checkmark & \texttimes \\
CTLE + K-means (ours)     & TinyStories-15M & 4    & 7.61\,MB        & $7.98\times$   & \phantom{0}19.23 & \textbf{MCU} & \checkmark & \texttimes \\
CTLE + GA (ours)          & TinyStories-15M & 4    & 7.61\,MB        & $7.98\times$   & \phantom{0}18.98 & \textbf{MCU} & \checkmark & \checkmark \\
\textbf{CTLE + PSO (ours)}& TinyStories-15M & \textbf{4} & \textbf{7.61\,MB} & $\mathbf{7.98\times}$ & \textbf{\phantom{0}15.93} & \textbf{MCU} & \checkmark & \checkmark \\
\bottomrule
\end{tabular}
\smallskip
\begin{minipage}{\textwidth}
  \footnotesize
  $^\dagger$ PPL reported for WikiText-2 on the referenced models; not comparable to our
  10-sentence TinyStories corpus.
  \textbf{Flash-stream}: inference reads weights directly from Flash without
  reconstructing dense matrices in RAM (LUT = 64\,B; only activations and KV cache allocated).
  GPU methods (GPTQ, SqueezeLLM, LLM.int8()) require VRAM for inference; CTLE runs on
  ESP32-P4 Nano (RISC-V HP\,@360\,MHz, 32\,MB PSRAM, 16\,MB Flash).
\end{minipage}
\end{table*}


\begin{table*}[t]
\centering
\caption{Comparison with representative post-training compression methods.
         \textbf{Note:} PPL values are evaluated on different models and datasets
         (see Model column) and are \emph{not} directly comparable.
         The table contrasts methodological characteristics and deployment target.
         $\ddagger$ FP32 footprint of counted tensors;
         $\star$ compression vs.\ FP32;
         $\dagger$ reported on WikiText-2, not our TinyStories corpus.}
\label{tab:sota}
\setlength{\tabcolsep}{4pt}
\renewcommand{\arraystretch}{1.2}
\begin{tabular}{llccccccc}
\toprule
\textbf{Method} & \textbf{Model} & \textbf{Bits/w} & \textbf{Size} & \textbf{Ratio$^\star$} &
\textbf{PPL\,$\downarrow$} & \textbf{HW} & \textbf{Flash-} & \textbf{Meta-} \\
 & & & & & & \textbf{target} & \textbf{stream} & \textbf{heuristic} \\
\midrule
FP32 (ours)               & TinyStories-15M & 32   & 60.8\,MB$^\ddagger$ & $1\times$      & \phantom{0}6.92  & MCU     & \texttimes & \texttimes \\
\midrule
GPTQ~\cite{frantar2023gptq}       & OPT-125M        & 4    & $\sim$15\,MB    & $\sim 8\times$ & 31.12$^\dagger$  & GPU     & \texttimes & \texttimes \\
SqueezeLLM~\cite{kim2024squeezellm}& LLaMA-7B        & 3.45 & $\sim$3.5\,GB   & $\sim 4\times$ & \phantom{0}5.79$^\dagger$ & GPU & \texttimes & \texttimes \\
LLM.int8()~\cite{dettmers2022llmint8}& OPT-$\geq$6.7B& 8    & $\sim 2\times$ red. & $\sim 2\times$ & $\approx$FP16$^\dagger$ & GPU & \texttimes & \texttimes \\
\midrule
INT4 Uniform (ours)       & TinyStories-15M & 4    & 7.61\,MB        & $7.99\times$   & 11070            & \textbf{MCU} & \checkmark & \texttimes \\
INT4 Block-wise (ours)    & TinyStories-15M & $\approx$5 & 9.51\,MB  & $6.39\times$   & \phantom{0}10.71 & \textbf{MCU} & \checkmark & \texttimes \\
CTLE + K-means (ours)     & TinyStories-15M & 4    & 7.61\,MB        & $7.98\times$   & \phantom{0}19.23 & \textbf{MCU} & \checkmark & \texttimes \\
CTLE + GA (ours)          & TinyStories-15M & 4    & 7.61\,MB        & $7.98\times$   & \phantom{0}18.98 & \textbf{MCU} & \checkmark & \checkmark \\
\textbf{CTLE + PSO (ours)}& TinyStories-15M & \textbf{4} & \textbf{7.61\,MB} & $\mathbf{7.98\times}$ & \textbf{\phantom{0}15.93} & \textbf{MCU} & \checkmark & \checkmark \\
\bottomrule
\end{tabular}
\smallskip
\begin{minipage}{\textwidth}
  \footnotesize
  $^\dagger$ PPL reported for WikiText-2 on the referenced models; not comparable to our
  10-sentence TinyStories corpus.
  \textbf{Flash-stream}: inference reads weights directly from Flash without
  reconstructing dense matrices in RAM (LUT = 64\,B; only activations and KV cache allocated).
  GPU methods (GPTQ, SqueezeLLM, LLM.int8()) require VRAM for inference; CTLE runs on
  ESP32-P4 Nano (RISC-V HP\,@360\,MHz, 32\,MB PSRAM, 16\,MB Flash).
\end{minipage}
\end{table*}


\begin{table*}[t]
\centering
\caption{Comparison with representative post-training compression methods.
         \textbf{Note:} PPL values are evaluated on different models and datasets
         (see Model column) and are \emph{not} directly comparable.
         The table contrasts methodological characteristics and deployment target.
         $\ddagger$ FP32 footprint of counted tensors;
         $\star$ compression vs.\ FP32;
         $\dagger$ reported on WikiText-2, not our TinyStories corpus.}
\label{tab:sota}
\setlength{\tabcolsep}{4pt}
\renewcommand{\arraystretch}{1.2}
\begin{tabular}{llccccccc}
\toprule
\textbf{Method} & \textbf{Model} & \textbf{Bits/w} & \textbf{Size} & \textbf{Ratio$^\star$} &
\textbf{PPL\,$\downarrow$} & \textbf{HW} & \textbf{Flash-} & \textbf{Meta-} \\
 & & & & & & \textbf{target} & \textbf{stream} & \textbf{heuristic} \\
\midrule
FP32 (ours)               & TinyStories-15M & 32   & 60.8\,MB$^\ddagger$ & $1\times$      & \phantom{0}6.92  & MCU     & \texttimes & \texttimes \\
\midrule
GPTQ~\cite{frantar2023gptq}       & OPT-125M        & 4    & $\sim$15\,MB    & $\sim 8\times$ & 31.12$^\dagger$  & GPU     & \texttimes & \texttimes \\
SqueezeLLM~\cite{kim2024squeezellm}& LLaMA-7B        & 3.45 & $\sim$3.5\,GB   & $\sim 4\times$ & \phantom{0}5.79$^\dagger$ & GPU & \texttimes & \texttimes \\
LLM.int8()~\cite{dettmers2022llmint8}& OPT-$\geq$6.7B& 8    & $\sim 2\times$ red. & $\sim 2\times$ & $\approx$FP16$^\dagger$ & GPU & \texttimes & \texttimes \\
\midrule
INT4 Uniform (ours)       & TinyStories-15M & 4    & 7.61\,MB        & $7.99\times$   & 11070            & \textbf{MCU} & \checkmark & \texttimes \\
INT4 Block-wise (ours)    & TinyStories-15M & $\approx$5 & 9.51\,MB  & $6.39\times$   & \phantom{0}10.71 & \textbf{MCU} & \checkmark & \texttimes \\
CTLE + K-means (ours)     & TinyStories-15M & 4    & 7.61\,MB        & $7.98\times$   & \phantom{0}19.23 & \textbf{MCU} & \checkmark & \texttimes \\
CTLE + GA (ours)          & TinyStories-15M & 4    & 7.61\,MB        & $7.98\times$   & \phantom{0}18.98 & \textbf{MCU} & \checkmark & \checkmark \\
\textbf{CTLE + PSO (ours)}& TinyStories-15M & \textbf{4} & \textbf{7.61\,MB} & $\mathbf{7.98\times}$ & \textbf{\phantom{0}15.93} & \textbf{MCU} & \checkmark & \checkmark \\
\bottomrule
\end{tabular}
\smallskip
\begin{minipage}{\textwidth}
  \footnotesize
  $^\dagger$ PPL reported for WikiText-2 on the referenced models; not comparable to our
  10-sentence TinyStories corpus.
  \textbf{Flash-stream}: inference reads weights directly from Flash without
  reconstructing dense matrices in RAM (LUT = 64\,B; only activations and KV cache allocated).
  GPU methods (GPTQ, SqueezeLLM, LLM.int8()) require VRAM for inference; CTLE runs on
  ESP32-P4 Nano (RISC-V HP\,@360\,MHz, 32\,MB PSRAM, 16\,MB Flash).
\end{minipage}
\end{table*}


\begin{table*}[t]
\centering
\caption{Comparison with representative post-training compression methods.
         \textbf{Note:} PPL values are evaluated on different models and datasets
         (see Model column) and are \emph{not} directly comparable.
         The table contrasts methodological characteristics and deployment target.
         $\ddagger$ FP32 footprint of counted tensors;
         $\star$ compression vs.\ FP32;
         $\dagger$ reported on WikiText-2, not our TinyStories corpus.}
\label{tab:sota}
\setlength{\tabcolsep}{4pt}
\renewcommand{\arraystretch}{1.2}
\begin{tabular}{llccccccc}
\toprule
\textbf{Method} & \textbf{Model} & \textbf{Bits/w} & \textbf{Size} & \textbf{Ratio$^\star$} &
\textbf{PPL\,$\downarrow$} & \textbf{HW} & \textbf{Flash-} & \textbf{Meta-} \\
 & & & & & & \textbf{target} & \textbf{stream} & \textbf{heuristic} \\
\midrule
FP32 (ours)               & TinyStories-15M & 32   & 60.8\,MB$^\ddagger$ & $1\times$      & \phantom{0}6.92  & MCU     & \texttimes & \texttimes \\
\midrule
GPTQ~\cite{frantar2023gptq}       & OPT-125M        & 4    & $\sim$15\,MB    & $\sim 8\times$ & 31.12$^\dagger$  & GPU     & \texttimes & \texttimes \\
SqueezeLLM~\cite{kim2024squeezellm}& LLaMA-7B        & 3.45 & $\sim$3.5\,GB   & $\sim 4\times$ & \phantom{0}5.79$^\dagger$ & GPU & \texttimes & \texttimes \\
LLM.int8()~\cite{dettmers2022llmint8}& OPT-$\geq$6.7B& 8    & $\sim 2\times$ red. & $\sim 2\times$ & $\approx$FP16$^\dagger$ & GPU & \texttimes & \texttimes \\
\midrule
INT4 Uniform (ours)       & TinyStories-15M & 4    & 7.61\,MB        & $7.99\times$   & 11070            & \textbf{MCU} & \checkmark & \texttimes \\
INT4 Block-wise (ours)    & TinyStories-15M & $\approx$5 & 9.51\,MB  & $6.39\times$   & \phantom{0}10.71 & \textbf{MCU} & \checkmark & \texttimes \\
CTLE + K-means (ours)     & TinyStories-15M & 4    & 7.61\,MB        & $7.98\times$   & \phantom{0}19.23 & \textbf{MCU} & \checkmark & \texttimes \\
CTLE + GA (ours)          & TinyStories-15M & 4    & 7.61\,MB        & $7.98\times$   & \phantom{0}18.98 & \textbf{MCU} & \checkmark & \checkmark \\
\textbf{CTLE + PSO (ours)}& TinyStories-15M & \textbf{4} & \textbf{7.61\,MB} & $\mathbf{7.98\times}$ & \textbf{\phantom{0}15.93} & \textbf{MCU} & \checkmark & \checkmark \\
\bottomrule
\end{tabular}
\smallskip
\begin{minipage}{\textwidth}
  \footnotesize
  $^\dagger$ PPL reported for WikiText-2 on the referenced models; not comparable to our
  10-sentence TinyStories corpus.
  \textbf{Flash-stream}: inference reads weights directly from Flash without
  reconstructing dense matrices in RAM (LUT = 64\,B; only activations and KV cache allocated).
  GPU methods (GPTQ, SqueezeLLM, LLM.int8()) require VRAM for inference; CTLE runs on
  ESP32-P4 Nano (RISC-V HP\,@360\,MHz, 32\,MB PSRAM, 16\,MB Flash).
\end{minipage}
\end{table*}
